\chapter{Корреляция, таблицы сопряжённости и дисперсионный анализ}
\label{ch:anova}

Главы~\ref{ch:basics}--\ref{ch:multiple} посвящены сравнению распределений и контролю ошибок при множественных тестах.
Настоящая глава переходит к структуре зависимостей между признаками и к моделям с факторами, разбивающими выборку на группы.
Если в главе~\ref{ch:hypothesis} мы спрашивали, отличаются ли средние, то здесь добавляется вопрос о силе связи признаков и о том, как несколько факторов совместно объясняют отклик.

В \S\ref{sec:dependencies} рассматриваются связи между непрерывными и категориальными переменными~--- корреляция Пирсона и Спирмена, критерий $\chi^2$ для таблиц сопряжённости.
В \S\ref{sec:anova} вводится дисперсионный анализ (ANOVA) \citep{fisher1925}~--- однофакторный, двухфакторный и модель с повторными измерениями (within-subjects).
Методы ANOVA тесно связаны с линейной регрессией (глава~\ref{ch:regression}). F-критерий для групповых средних~--- частный случай проверки коэффициентов в модели с dummy-переменными. Понимание разложения суммы квадратов облегчит переход к множественной регрессии.

\section{Зависимости между признаками}
\label{sec:dependencies}

\begin{definition}
Коэффициент корреляции Пирсона $r_{X_1X_2}$ измеряет силу линейной связи: \[ r_{X_1X_2}=\frac{\cov(X_1,X_2)}{\sqrt{\mathbb DX_1\,\mathbb DX_2}}.
\] Выборочный аналог $\hat r$ вычисляется по парам $(X_{1i},X_{2i})$.
\end{definition}

Корреляция отвечает на вопрос о совместной изменчивости. Если $X_1$ превышает своё среднее, насколько предсказуемо $X_2$ окажется больше или меньше своего.
Пирсон естественен при примерной нормальности и линейной связи. В главе~\ref{ch:basics} мы уже видели, что оценка моментов чувствительна к выбросам, и корреляция наследует эту чувствительность.

Критерий Стьюдента для $r$ (связанные выборки, $(X_1,X_2)\sim N(\mu,\Sigma)$): \[ H_0\colon \rho=0,\qquad T=\frac{\hat r\sqrt{n-2}}{\sqrt{1-\hat r^2}}\sim St(n-2).
\] Корреляция Спирмена~--- коэффициент Пирсона для рангов. Устойчива к монотонным нелинейным связям и выбросам.
Выборочная формула через ранги $X_{1i}$, $X_{2i}$: \[ \hat\rho_{X_1X_2}=1-\frac{6}{n^3-n}\sum_{i=1}^n \bigl(\rank(X_{1i})-\rank(X_{2i})\bigr)^2.
\] Критерий Стьюдента для $H_0\colon\rho=0$ имеет тот же вид, что и для Пирсона, с $\hat\rho$ вместо $\hat r$.
Когда связь монотонна, но нелинейна (например, экспоненциальная), Спирмен часто обнаруживает зависимость там, где $\hat r\approx 0$.

Корреляция Кендалла $\tau$~--- доля согласованных пар минус доля несогласованных: \[ \hat\tau_{X_1X_2}=\frac{C-D}{C+D}, \] где $C$~--- число пар $(i,j)$ с одинаковым порядком $(X_{1i},X_{2i})$ и $(X_{1j},X_{2j})$, $D$~--- с разным.
При $H_0$ и $n>10$ используют нормальную аппроксимацию. $\hat\tau$ обычно меньше $|\hat\rho|$ по модулю, но устойчивее к большим разрывам рангов и точнее на малых $n$.
При нормальности $\mathbb E\hat\tau=\mathbb E\hat\rho= \frac{2}{\pi}\arcsin r$ асимптотически.

Преобразование Фишера для построения доверительного интервала для $\rho$. $z=\operatorname{arctanh}(\hat r)$ асимптотически нормально с дисперсией $1/(n-3)$.
Корреляция не подразумевает причинно-следственной связи. Высокий $r$ может быть следствием общего скрытого фактора или совпадения.
Пирсон чувствителен к выбросам и обнаруживает преимущественно линейные зависимости. При U-образной связи $\hat r\approx 0$ при сильной функциональной зависимости.
Перед формальным тестом $H_0\colon\rho=0$ разумно построить диаграмму рассеяния. Глаз часто замечает нелинейность и выбросы раньше, чем один числовой коэффициент их скроет.
Две переменные с U-образной связью дадут $\hat r\approx 0$ и незначимый t-тест, тогда как Спирмена или квадратичная регрессия (гл.~\ref{ch:regression}) обнаружат зависимость.
Корреляция~--- не универсальный детектор связи, а мера линейной совместной изменчивости.
Выбор между Пирсоном, Спирменом и Кендаллом~--- не формальность. При тяжёлых хвостах и монотонной, но нелинейной связи ранговые меры дают более надёжный тест $H_0\colon\rho=0$. При линейной нормальной зависимости Пирсон максимально эффективен.
На примере температуры и широты маргинально широта сильно связана с температурой. После учёта долготы частная связь ослабевает или меняет знак~--- тот же механизм, что контроль ковариат в ANOVA и регрессии (гл.~\ref{ch:regression}).
Перед публикацией $r=0{,}85$ всегда стоит проверить, не смешивает ли третий фактор два признака.

\begin{definition}
Частная корреляция $r_{X_1X_2|X_3}$ измеряет линейную связь $X_1$ и $X_2$ после учёта влияния $X_3$: \[ r_{X_1X_2|X_3}=\frac{r_{X_1X_2}-r_{X_1X_3}r_{X_2X_3}} {\sqrt{(1-r_{X_1X_3}^2)(1-r_{X_2X_3}^2)}}.
\] Для нескольких управляющих признаков $M$ используют обратную матрицу выборочных корреляций $R=\Omega^{-1}$: $r_{X_iX_j|M\setminus\{X_i,X_j\}}=-R_{ij}/\sqrt{R_{ii}R_{jj}}$.
\end{definition}

Проверка $H_0\colon r_{X_1X_2|X_3}=0$ при $(X_1,X_2,X_3)\sim N(\mu,\Sigma)$: \[ T=\frac{\hat r_{X_1X_2|X_3}\sqrt{n-M-2}} {\sqrt{1-\hat r_{X_1X_2|X_3}^2}}\sim St(n-M-2), \] где $M$~--- число управляющих признаков.

Множественная корреляция $R_{X,M}$ оценивает, какую долю дисперсии $X$ объясняют признаки из $M$ линейно. $R_{X,M}\in[0,1]$.
Высокая маргинальная корреляция может исчезнуть или усилиться после учёта третьего фактора.
Температура января в городах США сильно коррелирует с широтой ($r\approx -0{,}85$), слабо~--- с долготой. Частная корреляция температуры с долготой при фиксированной широте значима ($r\approx 0{,}28$), отражая континентальный климат.
Этот пример показывает, почему простая корреляция редко бывает окончательным ответом. Смешивающие переменные могут как маскировать, так и создавать ложное впечатление связи.
Частная корреляция~--- первый шаг к контролю ковариат, которую в главе~\ref{ch:regression} мы реализуем через включение предикторов в регрессионную модель. Формула через обратную корреляционную матрицу предвосхищает стандартизованные коэффициенты множественной регрессии.
На рис.~\ref{fig:correlation} показана диаграмма рассеяния двух непрерывных признаков с линией регрессии. Визуальная оценка предшествует формальному тесту $H_0\colon\rho=0$.

\begin{figure}[htbp]
    \centering
    \includegraphics[width=\FigWidth]{\FigDir/chapter-4/corr_pearson.png}
    \caption{Диаграмма рассеяния и линейная регрессия.}
    \label{fig:correlation}
\end{figure}

Для таблицы $K_1\times K_2$ с частотами $n_{ij}$ и маргинальными суммами $n_{i+}$, $n_{+j}$, $n=\sum n_{ij}$: \[ H_0\colon \text{строка и столбец независимы},\qquad \chi^2=\sum_{i,j}\frac{\left(n_{ij}-\dfrac{n_{i+}n_{+j}}{n}\right)^2} {\dfrac{n_{i+}n_{+j}}{n}} =n\!\left(\sum_{i,j}\frac{n_{ij}^2}{n_{i+}n_{+j}}-1\right).
\] При $H_0$ и достаточных ожидаемых частотах ($\geq 5$ в большинстве ячеек) $\chi^2\sim\chi^2_{(K_1-1)(K_2-1)}$ асимптотически.
Условия~--- $n\geq 40$, ожидаемые частоты $<5$ не более чем в $20\%$ ячеек. Статистика измеряет отклонение наблюдаемых частот от ожидаемых при независимости. Большие значения указывают на систематическое совместное распределение категорий.

Точный критерий Фишера (таблица $2\times 2$). При фиксированных маргинальных суммах вероятность наблюдаемой таблицы \[ P=\frac{(a+b)!(c+d)!(a+c)!(b+d)!}{a!\,b!\,c!\,d!\,n!}.
\] p-value~--- сумма $P$ по всем таблицам с теми же суммами по строкам и столбцам, не более вероятных, чем наблюдаемая (двусторонний тест).
Для односторонней альтернативы $ad\ll bc$: \[ p=\sum_{i=0}^{a}\frac{C_{a+b}^{i}\,C_{c+d}^{a+c-i}}{C_n^{a+c}}.
\] На малых выборках асимптотический $\chi^2$ может быть ненадёжен. Точный критерий Фишера~--- стандартный выбор для таблиц $2\times 2$.
Для таблиц большего размера при нарушении условий на ожидаемые частоты используют точные перестановочные тесты или симуляцию. В прикладных пакетах (R, Python) выбор между асимптотическим и точным тестом обычно делается автоматически по правилам на минимальные ожидания.

\begin{definition}
Коэффициент $\phi$ (таблица $2\times 2$): $\phi=\sqrt{\chi^2/n}\in[-1,1]$.
Коэффициент V Крамера (общая $K_1\times K_2$): \[ V=\sqrt{\frac{\chi^2}{n\bigl(\min(K_1,K_2)-1\bigr)}}\in[0,1].
\] Нулевое значение~--- отсутствие связи. $V=1$~--- полное совпадение категорий (с точностью до переименования уровней).
\end{definition}

Корреляция Мэтьюса (бинарные признаки): $MCC=(ad-bc)/\sqrt{(a+b)(a+c)(b+d)(c+d)}$.
Значимый $\chi^2$ при большом $n$ не означает сильную связь. Коэффициенты $\phi$, $V$ и MCC измеряют силу ассоциации и должны сообщаться вместе с p-value, по аналогии с тем, как в главе~\ref{ch:hypothesis} важны не только значимость, но и величина эффекта.

Пример. Таблица образование~--- вера задаёт сопряжённость уровня образования (строки) и ответа на вопрос о религии (столбцы).
Критерий $\chi^2$ даёт $\chi^2\approx 76{,}1$, p-value $\ll 0{,}001$, степени свободы $(3-1)(6-1)=10$.

Распределение ответов зависит от образования. Например, доля ответа know God exists снижается с ростом образовательного уровня.
Остатки $(O-E)/\sqrt E$ показывают, какие ячейки вносят наибольший вклад в статистику.
Анализ остатков превращает глобальное отклонение от независимости в локальные содержательные утверждения о конкретных подгруппах.

Парадокс Симпсона. Разделим данные о лечении препаратом и плацебо по полу.
У мужчин препарат выигрывает ($\chi^2\approx 5{,}5$, p-value $\approx 0{,}02$), у женщин~--- тоже ($\chi^2\approx 17{,}6$, p-value $\ll 0{,}001$), но в объединённой таблице связи нет ($\chi^2\approx 0{,}4$, p-value $\approx 0{,}54$).
Препарат чаще давали мужчинам с тяжёлым течением, плацебо~--- женщинам с лёгким. Смешение пропорций переворачивает вывод.

Классический пример Беркли (1973). На уровне университета в целом мужчины поступали чаще женщин ($44{,}3\%$ против $34{,}6\%$, $\chi^2\approx 108$), но на каждом из крупнейших факультетов доля поступивших женщин не уступала мужчинам~--- женщины чаще подавали документы на конкурентные направления.
Содержательный вывод возможен только при согласованном разбиении или явном моделировании смешивающих факторов.
Парадокс Симпсона напоминает о роли частной корреляции. Маргинальная ассоциация может исчезнуть или сменить знак после учёта смешивающей переменной.
В прикладных исследованиях (медицина, социология, A/B-тестирование) таблицы сопряжённости часто сопровождают стратификацией по полу, возрасту или каналу трафика~--- именно чтобы не попасть в ловушку агрегированного вывода.
Пример лечения по полу с объединённой таблицей повторяет логику парадокса Симпсона на языке $\chi^2$. Значимость в подгруппах и отсутствие связи в агрегате~--- не противоречие формул, а следствие разных смешивающих пропорций.
Практический вывод~--- сначала стратификация по содержательным факторам, затем маргинальный тест, только если граф или предметная область оправдывает объединение.

При интерпретации таблиц сопряжённости следует избегать типичных ошибок.
\begin{itemize}
    \item Интерпретировать значимый $\chi^2$ как сильную связь.
          Для больших $n$ слабые отклонения от независимости значимы. Коэффициенты $\phi$, Cramér's $V$ измеряют силу связи.
    \item Применять асимптотический $\chi^2$ при малых ожидаемых
          частотах ($<5$) без точного критерия Фишера.
    \item Делать причинные выводы из таблицы сопряжённости
          без экспериментального дизайна.
\end{itemize}

\section{Дисперсионный анализ}
\label{sec:anova}

ANOVA разбивает общую вариабельность отклика на компоненты, обусловленные факторами и остаточную, и проверяет гипотезы о равенстве средних в группах при предположении нормальности и равенства дисперсий внутри групп.
Идея близка к $t$-критерию из главы~\ref{ch:hypothesis}. При двух группах однофакторный ANOVA эквивалентен двухвыборочному $t$-критерию, но при $K>2$ нужен единый общий критерий, прежде чем локализовать различия между парами.

Пусть $K$ групп, $n_k$ наблюдений в группе $k$, всего $N=\sum n_k$.
Разложение суммы квадратов: \[ \sum_{k,i}(Y_{ki}-\bar Y)^2= \underbrace{\sum_{k,i}(Y_{ki}-\bar Y_k)^2}_{\mathrm{SSW}}+ \underbrace{\sum_k n_k(\bar Y_k-\bar Y)^2}_{\mathrm{SSB}}, \] $\mathrm{SST}=\mathrm{SSW}+\mathrm{SSB}$.
Средние квадраты: $\mathrm{MSB}=\mathrm{SSB}/(K-1)$, $\mathrm{MSW}=\mathrm{SSW}/(N-K)$.
\[ H_0\colon \mu_1=\cdots=\mu_K,\qquad F=\frac{\mathrm{MSB}}{\mathrm{MSW}}\sim F(K-1,N-K)\ \text{при}\ H_0.
\] Межгрупповая сумма квадратов $\mathrm{SSB}$ отражает, насколько средние групп разбросаны вокруг общего среднего. Внутригрупповая $\mathrm{SSW}$~--- шум внутри групп. Большой $F$ означает, что междугрупповая вариация доминирует над внутригрупповой.
Предпосылки ANOVA~--- нормальность внутри групп и равенство дисперсий (гомоскедастичность)~--- проверяют до интерпретации. При их нарушении F-критерий может сохранять уровень значимости при больших $N$, но мощность и интервалы для разностей средних страдают.
Перед ANOVA имеет смысл построить ящики с усами по группам и проверить равенство дисперсий (Левена, F на $\sigma^2$). Нарушение гомоскедастичности при разных $n_k$ особенно опасно. F-тест может оставаться значимым, хотя различие средних объясняется смешением разброса и центра.
Робастная альтернатива~--- Крускала--Уоллиса с последующим апостериорным анализом по рангам.

\begin{definition}
Eta-squared (доля объяснённой дисперсии в выборке): $\eta^2=\mathrm{SSB}/\mathrm{SST}$.
Omega-squared (несмещённая оценка в популяции): \[ \hat\omega^2=\frac{\mathrm{SSB}-\mathrm{MSW}(K-1)} {\mathrm{SST}+\mathrm{MSW}/(N-K)}.
\] При отвержении $H_0$ важно сообщать не только p-value, но и $\eta^2$ или $\hat\omega^2$. При больших $N$ слабый эффект может быть статистически значимым.
\end{definition}

Критерий Крускала--Уоллиса (непараметрический общий). Ранжируем все $N$ наблюдений, $H_0\colon$ распределения в группах совпадают с точностью до сдвига. Без связок: \[ K=\frac{12}{N(N+1)}\sum_{k=1}^K n_k\bar r_k^2-3(N+1) \;\dot\sim\;\chi^2_{K-1}\quad (n_k>5).
\] HSD Тьюки (апостериорный, контролирует FWER попарных сравнений): \[ \mathrm{HSD}=\frac{q_\alpha(N-K)\,S}{\sqrt{\bar n}}, \quad S^2=\frac{\mathrm{SSW}}{N-K}, \quad \bar n=\frac{K}{\sum_k 1/n_k}.
\] Если $|\bar Y_i-\bar Y_j|>\mathrm{HSD}$, то $H_0\colon\mu_i=\mu_j$ отвергается. Ранговый аналог~--- критерий Неменьи.
При нарушении нормальности при больших $n$ F-критерий устойчив (robust). При неравных дисперсиях и разных $n_k$ предпочтителен критерий Уэлча или непараметрический Крускала--Уоллиса.
ANOVA проверяет наличие различий. Апостериорные тесты (HSD Тьюки) локализуют пары групп.
Связь с главой~\ref{ch:multiple} прямая. При $K$ группах попарных сравнений $K(K-1)/2$, и HSD Тьюки~--- один из способов контроля FWER без процедуры всех попарных $t$-тестов подряд.

Пример. В классическом примере ToothGrowth ($N=60$) измерена длина роста зубов (\texttt{len}) у крыс при двух типах добавки (\texttt{supp}: OJ, VC) и трёх дозах витамина~C (\texttt{dose}: 0.5, 1.0, 2.0).

На рис.~\ref{fig:anova-box} показаны ящики с усами по комбинациям факторов.
Однофакторный ANOVA по \texttt{supp} выявляет различие средних (p-value $\ll 0{,}001$). Апельсиновый сок (OJ) в среднем эффективнее аскорбата (VC).
По \texttt{dose} эффект дозы также значим. Рост увеличивается с дозой.
После значимого общего критерия HSD Тьюки локализует различия. Например, OJ превосходит VC, доза $2{,}0$~--- $0{,}5$ и $1{,}0$.
Ящики с усами помогают до формального теста увидеть, совпадают ли дисперсии и нет ли выбросов, нарушающих предпосылки ANOVA.
Класс ToothGrowth часто используют в учебных целях именно потому, что факторы \texttt{supp} и \texttt{dose} наглядны, а число наблюдений достаточно для устойчивого F-теста, но невелико~--- студент может наглядно увидеть различие между общим критерием и апостериорной локализацией.
Однофакторный ANOVA по \texttt{dose} отвечает на вопрос, есть ли эффект дозы вообще. HSD Тьюки отвечает, какие пары доз различаются.
Значимый общий критерий без апостериорного анализа лишь констатирует наличие различий.
Значимый апостериорный анализ при незначимом общем критерии при $K=2$ эквивалентен одному $t$-тесту. При $K>2$ он невозможен без предварительного F-теста (контроль FWER).

\begin{figure}[htbp]
    \centering
    \includegraphics[width=\FigWidth]{\FigDir/chapter-4/drugs.png}
    \caption{Рост зубов по типу добавки и дозе витамина~C.}
    \label{fig:anova-box}
\end{figure}

Для факторов $A$ ($I$ уровней) и $B$ ($J$ уровней), по $n$ наблюдений в ячейке: \[ Y_{ijk}=\mu+\alpha_i+\beta_j+(\alpha\beta)_{ij}+\varepsilon_{ijk}, \quad \varepsilon_{ijk}\sim N(0,\sigma^2).
\] Разложение дисперсии (сбалансированный дизайн): \[ \mathrm{SS}_{total}=\mathrm{SS}_A+\mathrm{SS}_B+\mathrm{SS}_{AB}+\mathrm{SS}_{res}, \] \[ \mathrm{SS}_A=nJ\sum_i(\bar Y_{i\bullet}-\bar Y)^2,\quad \mathrm{SS}_B=nI\sum_j(\bar Y_{\bullet j}-\bar Y)^2, \] \[ \mathrm{SS}_{AB}=n\sum_{i,j} (\bar Y_{ij}-\bar Y_{i\bullet}-\bar Y_{\bullet j}+\bar Y)^2.
\] F-статистики: $F_A=\mathrm{MS}_A/\mathrm{MS}_{res}$, $F_B=\mathrm{MS}_B/\mathrm{MS}_{res}$, $F_{AB}=\mathrm{MS}_{AB}/\mathrm{MS}_{res}$.
Член $(\alpha\beta)_{ij}$ описывает взаимодействие. Эффект одного фактора зависит от уровня другого.
Интерпретировать main effects при значимом $A\times B$ без осторожности нельзя. Средний эффект фактора может не представлять ни одной реальной комбинации уровней.
Двухфакторный дизайн ToothGrowth естественно расширяется до модели с \texttt{supp}, \texttt{dose} и их произведением. Именно взаимодействие может объяснить, почему OJ эффективнее VC на одних дозах и нет на других, если такая картина видна на рис.~\ref{fig:anova-box}.

Пример. В данных об артериальном давлении измерено артериальное давление (\texttt{pressure}) при комбинациях биологической обратной связи (\texttt{biofeedback}: present/absent), диеты (\texttt{diet}) и препарата (\texttt{drug}: 1, 2, 3).

Двухфакторный ANOVA (biofeedback $\times$ diet) внутри \texttt{drug=1} даёт значимый эффект biofeedback ($F\approx 5{,}9$, p-value $\approx 0{,}024$). Присутствие обратной связи снижает давление.
Для \texttt{drug=2} эффект не значим (p-value $\approx 0{,}44$), что иллюстрирует взаимодействие препарата с не медицинским вмешательством.
Раздельный анализ по уровням \texttt{drug} здесь не прихоть, а содержательная необходимость. Объединённый ANOVA смешал бы разные механизмы действия препаратов.
Интуиция взаимодействия $A\times B$. Эффект фактора $A$ измеряется по-разному на разных уровнях $B$.
На графике \texttt{supp} $\times$ \texttt{dose} линии для OJ и VC могут пересекаться~--- тогда main effect \texttt{supp} без учёта взаимодействия вводит в заблуждение. Утверждение, что в среднем OJ лучше, не описывает ни одну конкретную дозу.
Такой график~--- обязательное дополнение к таблице F-тестов.

Within-subjects (повторные измерения). Один и тот же объект наблюдается во всех уровнях фактора $A$.
Модель: \[ Y_{ij}=\mu+\alpha_i+\gamma_j+\varepsilon_{ij}, \] где $\gamma_j$~--- эффект объекта $j$, $\sum\alpha_i=0$.
Разложение: \[ \mathrm{SS}_{total}=\mathrm{SS}_A+\mathrm{SS}_{subj}+\mathrm{SS}_{err}, \] \[ F=\frac{\mathrm{SS}_A/(K-1)} {(\mathrm{SS}_{wg}-\mathrm{SS}_{subj})/\bigl((n-1)(K-1)\bigr)} \sim F\bigl(K-1,(n-1)(K-1)\bigr)\ \text{при}\ H_0, \] где $\mathrm{SS}_{wg}$~--- общая внутрисубъектная вариабельность, $\mathrm{SS}_{subj}=K\sum_j(\bar Y_{\cdot j}-\bar Y)^2$~--- между субъектами.
Сферичность (Mauchly). Для $K>2$ все попарные разности $Y_{ik}-Y_{il}$ имеют одинаковую дисперсию при любом $i$.
Нарушение сферичности завышает вероятность ошибки I рода.
Поправки. Greenhouse--Geisser и Huynh--Feldt умножают числитель и знаменатель $F$ на коэффициент $\varepsilon\in(0,1]$, оцениваемый по ковариационной матрице. $\varepsilon=1$ при сферичности.
При сильном нарушении используют критерий Фридмана (ранговый общий) или multivariate подход.
Повторные измерения повышают мощность за счёт контроля межиндивидуальной вариабельности, но требуют выполнения сферичности ковариационной матрицы (критерий Mauchly).
При нарушении сферичности применяют поправки Greenhouse--Geisser или Huynh--Feldt. При $K=2$ сферичность выполняется автоматически.
Содержательно мы сравниваем изменение отклика у одних и тех же субъектов, а не смешиваем межиндивидуальные различия с эффектом условия.
Repeated-measures ANOVA~--- частный случай смешанной модели (глава~\ref{ch:regression}). Фиксированный эффект условия и случайный эффект субъекта разделяют два источника вариабельности, что особенно ценно, когда межсубъектный разброс велик (как в экспериментах с временем реакции).

Пример. Таблица данных о времени реакции на каннабис содержит время реакции ($\mu$с) для $12$ испытуемых при трёх условиях. Без воздействия (\texttt{None}), лёгком (\texttt{Light}) и умеренном (\texttt{Moderate}) употреблении марихуаны (столбцы). Строки~--- субъекты.

Это классическая схема repeated measures. Каждый испытуемый измерен трижды.
Однофакторный ANOVA с фактором условие и блоком субъект позволяет проверить $H_0\colon$ среднее время реакции не зависит от дозы, контролируя индивидуальные различия.
Профильные графики показывают рост времени реакции с дозой у большинства субъектов. При $K=3$ условиях имеет смысл проверить сферичность и при необходимости применить поправку G--G.
Ошибка~--- трактовать $36=12\times 3$ наблюдений как независимые. Это завышает мощность и искажает p-value, поскольку три измерения одного человека не независимы.
Парный подход из главы~\ref{ch:hypothesis} (критерий для разностей после минус до внутри субъекта) даёт эквивалент repeated-measures выводу при двух условиях. При $K>2$ нужна полная ANOVA-модель с блоком субъект.
Профиль до, лёгкое и умеренное у одного испытуемого~--- связанные наблюдения. Игнорирование субъекта эквивалентно завышению $n$ и занижению p-value.
Содержательно мы спрашиваем, меняется ли реакция внутри человека, а не отличаются ли 36 независимых измерений.
Repeated-measures ANOVA формализует этот вопрос через блок субъект и F-тест на эффект условия.

При планировании ANOVA следует помнить о типичных ловушках.
\begin{itemize}
    \item Применять однофакторный ANOVA к времени реакции без учёта
          субъекта, трактуя $36$ наблюдений как независимые~--- завышает мощность и искажает p-value.
    \item Игнорировать взаимодействие в гипертонии и
          интерпретировать main effects при значимом $A\times B$.
    \item Считать незначимый F-тест доказательством отсутствия эффекта
          без анализа мощности и доверительных интервалов для разностей.
\end{itemize}

Исследование связи между часами обучения ($X$) и баллом экзамена ($Y$) для $n=200$ студентов.
Диаграмма рассеяния показывает линейный тренд, $\hat r=0{,}52$, $t$-тест на $\rho=0$ даёт p-value $\ll 0{,}001$.
Известен фактор факультет (гуманитарный / технический). Частная корреляция $r_{XY|Z}$ падает до $0{,}21$~--- часть связи объясняется разным распределением часов по факультетам.

Параллельно однофакторный ANOVA балл по факультету ($K=3$ группы) даёт значимый $F$, $\eta^2=0{,}14$.
Апостериорный HSD локализует различие между гуманитарным и техническим.
Маргинальная корреляция часов и баллов завышена из-за смешивания групп. ANOVA и частная корреляция согласуются в том, что факультет~--- смешивающий фактор в регрессионном смысле (гл.~\ref{ch:regression}).

Если бы мы игнорировали $Z$ и строили только регрессию часов на балл, политика больше часов всем могла бы не работать внутри факультета.
Такой пример связывает \S~\ref{sec:dependencies} и \S~\ref{sec:anova}. Разные инструменты, один вопрос о структуре зависимостей.
График (диаграмму рассеяния с цветом по факультету) следует построить до любого из формальных тестов.

Прикладной отчёт часто смешивает типы переменных. Непрерывные предикторы (корреляция, регрессия), категориальные (таблицы $\chi^2$), групповые средние (ANOVA).
Ошибка~--- применить один тест по привычке. $\chi^2$ к дискретизированному непрерывному признаку без обоснования интервалов. ANOVA к парным данным без блока субъект. Корреляцию Пирсона к монотонной нелинейной связи.
Явная строка тип переменной и метод в начале анализа предотвращает такие подмены и связывает \S~\ref{sec:dependencies} и \S~\ref{sec:anova} в единую методологию.

Dummy-кодирование фактора с $K$ уровнями и однофакторный ANOVA эквивалентны линейной модели с $K-1$ индикаторами.
F-тест на все dummy~--- общий ANOVA. t-тест на один коэффициент~--- contrast.
Понимание ANOVA через регрессию (гл.~\ref{ch:regression}) снимает ощущение другого предмета. $\mathrm{SSB}$~--- объяснённая регрессией доля $TSS$ при групповых предикторах.
Two-way ANOVA с interaction соответствует модели с произведениями dummy.

Значимый общий F или $\chi^2$ при большом $n$ совместим с малым $\eta^2$ или Cramér's $V$.
В отчёте указывают p-value, меру силы ($\eta^2$, $\phi$, $d$ Коэна для контрастов), график (ящики, мозаичная диаграмма).
Апостериорный анализ только после значимого общего критерия (кроме запланированных контрастов).
Repeated measures требуют явно указать субъект в модели.

\section{G-критерий, перестановочная корреляция и робастный ANOVA}

G-критерий для таблиц сопряжённости. Логарифмическое отношение правдоподобия для независимости в таблице $K_1\times K_2$: \[ G^2 = 2\sum_{i,j} n_{ij}\ln\frac{n_{ij}n}{n_{i+}n_{+j}} = 2n\!\left(\sum_{i,j}\frac{n_{ij}}{n}\ln\frac{n_{ij}n}{n_{i+}n_{+j}}\right).
\] При $H_0$ и больших ожидаемых частотах $G^2\dot\sim\chi^2_{(K_1-1)(K_2-1)}$, как и статистика Пирсона $\chi^2$. Численно $G^2\approx\chi^2$.
G-критерий удобен в логлинейных и GLM-моделях (гл.~\ref{ch:regression}).

\begin{critcard}{G-критерий независимости}
Выборка & связанные $(X_{1i},X_{2i})$, $X_1\in\{1,\ldots,K_1\}$, $X_2\in\{1,\ldots,K_2\}$ \\
$H_0$ & $X_1$ и $X_2$ независимы \\
Статистика & $G^2 = 2\sum_{i,j} n_{ij}\ln\bigl(n_{ij}n/(n_{i+}n_{+j})\bigr)$ \\
Нулевое распр. & $\chi^2_{(K_1-1)(K_2-1)}$ (асимптотика) \\
p-value & хвост $\chi^2$; те же условия на ожидаемые частоты, что у $\chi^2$ Пирсона \\
\end{critcard}

Перестановочный критерий для корреляции $r$. Для связанных выборок $(X_{1i},X_{2i})$ при $H_0\colon\rho=0$ переставляют индексы одной из выборок ($n!$ вариантов).
На каждой перестановке $g$ вычисляют $\hat r(X_1^n, gX_2^n)$. p-value~--- доля перестановок с $|\hat r|\geq|\hat r_{\text{obs}}|$ (двусторонний тест).

\begin{critcard}{Перестановочный тест $H_0\colon\rho=0$}
Выборки & связанные $X_1^n$, $X_2^n$ \\
$H_0$ & $\rho=0$ (линейная корреляция Пирсона) \\
Статистика & $T=\hat r_{X_1X_2}$ \\
Нулевое распр. & эмпирическое по перестановкам $X_2^n$ \\
p-value & $\#\{g\colon|\hat r_g|\geq|\hat r_{\text{obs}}|\}/n!$ (или Монте-Карло) \\
\end{critcard}

Критерий Джонкхеера--Терпстра. Omnibus для упорядоченных групп. $H_0\colon F_1=\cdots=F_K$ против $H_1\colon F_1\preceq\cdots\preceq F_K$ (стохастический порядок).
Статистика $S=\sum_{k,i} a_{ki}$, где $a_{ki}$~--- число наблюдений в группах $1,\ldots,k-1$, меньших $X_{ki}$.
При $n_k>10$: $S\sim N(\mu,\sigma^2)$ с формулами $\mu,\sigma$ из лекции~6.

\begin{critcard}{Джонкхеер--Терпстра (тренд по группам)}
Выборки & $K$ независимых групп, размеры $n_k$ \\
$H_0$ & $\med X_1=\cdots=\med X_K$ (равные распределения) \\
$H_1$ & $\med X_1\leq\cdots\leq\med X_K$ (не все равны) \\
Статистика & $S=\sum_{k,i}a_{ki}$ (число меньших слева) \\
Нулевое распр. & табличное; нормальная аппроксимация при больших $n_k$ \\
\end{critcard}

Критерий Даннета. Сравнение $K-1$ лечений с одной контрольной группой (индекс~1): \[ D_i = \frac{\bar X_i-\bar X_1}{S\sqrt{1/n_i+1/n_1}},\quad S^2=\frac{1}{N-K}\sum_k (n_k-1)S_k^2.
\] При нормальности и $\mu_1=\cdots=\mu_K$ вектор $(D_2,\ldots,D_K)$ имеет многомерное $t$-распределение. Пороги дают FWER для семейства каждая группа против контроля.

\begin{critcard}{Даннет (все группы против контроля)}
Выборки & $K$ нормальных групп, общая $\sigma^2$ \\
$H_0$ & $\mu_1=\cdots=\mu_K$ \\
Частные $H_0^{(i)}$ & $\mu_i=\mu_1$, $i=2,\ldots,K$ \\
Статистика & $D_i$ как выше \\
Нулевое распр. & многомерное $t$ (таблицы Даннета) \\
Гарантия & контроль FWER для семейства сравнений с контролем \\
\end{critcard}

LSD Фишера. После значимого общего ANOVA попарно сравнивают \[ \frac{\bar X_i-\bar X_j}{S\sqrt{1/n_i+1/n_j}}\sim St(n_i+n_j-2).
\] Наименьшая значимая разность: \[ \mathrm{LSD}_{ij}=\frac{t_{\alpha,\,n_i+n_j-2}\,S}{\sqrt{1/n_i+1/n_j}}.
\] Если $|\bar X_i-\bar X_j|>\mathrm{LSD}_{ij}$, отвергают $H_0\colon\mu_i=\mu_j$.

\begin{remark}
LSD не контролирует FWER при множественных попарных сравнениях. Использовать только как исследовательский инструмент после $F$-теста, не как замену HSD Тьюки или Холма.
\end{remark}

Критерий Бартлетта (однородность дисперсий). Проверка $H_0\colon\sigma_1^2=\cdots=\sigma_K^2$ перед ANOVA: \[ B = \frac{\ln 10}{C}\left((N-K)\ln S^2 - \sum_k (n_k-1)\ln S_k^2\right), \quad C = 1+\frac{1}{3(K-1)}\left(\sum_k\frac{1}{n_k-1}-\frac{1}{N-K}\right), \] $S^2$~--- объединённая дисперсия. При $n_k>6$: $B\dot\sim\chi^2_{K-1}$.
Критерий чувствителен к ненормальности. При сомнениях~--- Левена.

\begin{critcard}{Бартлетт (гомогенность $\sigma^2$)}
Выборки & $X_{ki}\sim N(\mu_k,\sigma_k^2)$ \\
$H_0$ & $\sigma_1^2=\cdots=\sigma_K^2$ \\
Статистика & $B(X^N)$ как выше \\
Нулевое распр. & $\chi^2_{K-1}$ (аппроксимация) \\
\end{critcard}

Критерий Неменьи (CD). Ранговый апостериорный анализ после Крускала--Уоллиса. Критическая разность: \[ \mathrm{CD} = q'_\alpha\sqrt{\frac{K(K+1)}{6N}}, \] где $q'_\alpha$~--- квантиль стьюдентизированного размаха для рангов.
Если $|\bar r_i-\bar r_j|>\mathrm{CD}$, отвергают $H_0\colon$ сдвиг $i$-й и $j$-й групп совпадает.

\begin{critcard}{Неменьи (попарно по рангам)}
Предпосылка & значимый или плановый Крускала--Уоллиса \\
$H_0^{(ij)}$ & распределения в группах $i,j$ совпадают (по рангам) \\
Статистика & $|\bar r_i-\bar r_j|$ и $\mathrm{CD}$ \\
Гарантия & FWER для попарных сравнений в ранговой постановке \\
\end{critcard}

Критерии Фридмана и Пейджа (повторные измерения). Для $n$ субъектов и $K$ условий ранжируют отклик внутри каждой строки (субъекта). $r_{ki}$~--- ранг $k$-го условия у $i$-го субъекта, $R_k=\sum_i r_{ki}$.
\[ S = \frac{12}{nK(K+1)}\sum_{k=1}^K R_k^2 - 3n(K+1) \quad\text{(Фридман, общий)}.
\] Пейдж (тренд): $L=\sum_{k=1}^K k R_k$. При $H_0$ и больших $n,K$: $L\sim N(\mu_L,\sigma_L^2)$ с параметрами из лекции~6.

\begin{critcard}{Фридман (within-subjects, общий)}
Данные & матрица $X_{ki}$, $i=1,\ldots,n$, $k=1,\ldots,K$ \\
$H_0$ & $\alpha_1=\cdots=\alpha_K$ (эффект условия отсутствует) \\
Статистика & $S$ через ранги $r_{ki}$ в строках \\
Нулевое распр. & $\chi^2_{K-1}$ при $n>15$, $K>10$ \\
\end{critcard}

\begin{critcard}{Пейдж (тренд условий)}
$H_0$ & как у Фридмана \\
$H_1$ & $\alpha_1\leq\cdots\leq\alpha_K$ (монотонный тренд) \\
Статистика & $L=\sum_k k R_k$ \\
Нулевое распр. & $N(\mu_L,\sigma_L^2)$ (аппроксимация) \\
\end{critcard}

Критерий Крускала--Уоллиса (кратко). Ранговый общий для $K$ независимых групп (дополнение к \S~\ref{sec:anova}). Ранжировать все $N$ наблюдений, $\bar r_k$~--- средний ранг в группе $k$, \[ H = \frac{12}{N(N+1)}\sum_{k=1}^K n_k\bar r_k^2 - 3(N+1) \;\dot\sim\;\chi^2_{K-1}\quad (n_k>5).
\] После значимого $H$~--- апостериорный Неменьи. При упорядоченных уровнях фактора~--- Джонкхеер вместо $H$.

\begin{critcard}{Крускал--Уоллис}
Выборки & $K$ независимых групп, $N=\sum n_k$ \\
$H_0$ & распределения совпадают (сдвиг не отличается) \\
Статистика & $H$ (или $K$ в обозначениях лекции) через $\bar r_k$ \\
Нулевое распр. & $\chi^2_{K-1}$ \\
\end{critcard}

При сильной асимметрии или выбросах Крускала--Уоллиса заменяет F-тест. Попарные сравнения~--- по рангам (Немени).
Дисперсионный анализ Уэлча обобщает $t$-Уэлча на $K>2$ групп при неравных дисперсиях.
Выбор процедуры следует той же схеме, что в гл.~\ref{ch:hypothesis}. График, предпосылки, общий критерий, апостериорный анализ с поправкой.

\section{Практика и семинар}
\label{sec:practice-ch4}

Задания соответствуют семинарам~5--6. Для каждой задачи зафиксируйте тип переменных, $H_0/H_1$, предпосылки, числовой вывод (p-value, размер эффекта) и содержательную интерпретацию в 3--5 предложениях.

Корреляция и выброс (sem5). Модель $X_1\sim N(0,1)$, $X_2=X_1+\varepsilon$, $\varepsilon\sim N(0,\alpha^2)$, $n=10$ (затем $n=10^5$). Теоретически $r=1/\sqrt{1+\alpha^2}$.

Шаги.

\begin{enumerate}
    \item Сгенерировать пары. Вычислить $\hat r$, $t$-тест $H_0\colon\rho=0$
          и ДИ через преобразование Фишера. Сравнить с теорией.
    \item Для $\alpha\in\{0{,}1,0{,}5,\ldots,5\}$ (шум) в цикле $1000$ раз
          оценить долю принятий $H_0$ при истинном $r\neq 0$ (мощность/ошибка II рода для малых $n$).
    \item Выброс. При $n=10^5$, $\alpha=1$ добавить к одной точке
          сдвиг $+9999$ по $X_1$. Сравнить $\hat r$ Пирсона до/после. Повторить для Спирмена и Кендалла.
    \item Перестановочный тест. Для данных без выброса оценить
          $p_{\text{perm}}$ (перестановка $X_2$). Сравнить с $p_t$.
    \item Интерпретация. Объяснить, почему один выброс меняет
          вывод Пирсона, но почти не меняет ранговые коэффициенты. Когда перестановочный p-value заметно отличается от $t$.
\end{enumerate}

Гипертония. Трёхфакторный ANOVA (sem6). Данные.\texttt{biofeedback} $\times$ \texttt{diet} $\times$ \texttt{drug} $\to$ \texttt{pressure} ($N=72$, Maxwell \& Delaney).

Шаги.

\begin{enumerate}
    \item Построить ящики с усами по \texttt{drug} и фасетам
          \texttt{biofeedback}$\times$\texttt{diet}.
    \item Оценить трёхфакторную модель с взаимодействиями
          \texttt{drug*diet*biofeedback} (Type~II/III ANOVA).
    \item Если $A\times B\times C$ значимо~--- не интерпретировать
          main effects без графика. Разобрать простые эффекты (например, \texttt{biofeedback} внутри \texttt{drug=1}).
    \item Сопоставить с примером \S~\ref{sec:anova}. Почему отдельный ANOVA
          по уровням \texttt{drug} содержательно оправдан.
    \item Интерпретация. Сформулировать, какие комбинации лечения
          снижают давление. Указать $F$, p-value, $\eta^2$ для значимых членов.
\end{enumerate}

Хоровые певцы. Бартlett и Тьюки (sem6). Данные. Рост (дюймы) и регистр голоса (\texttt{voice}: bass, baritone, tenor, alto), $n=235$.

Шаги.

\begin{enumerate}
    \item Визуализация (ящики, violin plot).
    \item \texttt{scipy.stats.bartlett} по группам. $H_0\colon\sigma^2$
          одинаковы. При p-value $<0{,}05$~--- дисперсионный анализ Уэлча или Крускала--Уоллиса вместо классического $F$.
    \item Однофакторный ANOVA (или робастный аналог). При значимом общем критерии~---
          \texttt{pairwise\_tukeyhsd} / \texttt{multicomp}.
    \item Интерпретация. Есть ли систематическая связь роста
          с регистром? Указать, какие пары различаются по апостериорному анализу. Обсудить нарушение гомогенности дисперсий, если оно есть.
\end{enumerate}

Витамин~C и рост зубов морских свинок (sem6). Данные.\texttt{ToothGrowth}: \texttt{len}, \texttt{supp} (OJ/VC), \texttt{dose} (0.5, 1.0, 2.0).

Шаги.

\begin{enumerate}
    \item График \texttt{len} и \texttt{dose}, цвет по \texttt{supp}.
    \item Двухфакторный ANOVA \texttt{len ~ supp * dose}. Проверить
          взаимодействие на графике (параллельность линий).
    \item При значимом общем критерии по \texttt{dose}~--- HSD Тьюки для доз.
          Отдельно прокомментировать \texttt{supp}.
    \item Сравнить с табл.~лекции (HSD для жаропонижающих). Логика та же.
    \item Интерпретация. Какая добавка и доза максимальны.
          Есть ли пересечение кривых OJ/VC?
\end{enumerate}

Марихуана и время реакции (sem6). Данные $12$ субъектов $\times$ $3$ условия (\texttt{None}, \texttt{Light}, \texttt{Moderate}). Отклик~--- время реакции ($\mu$с).

Шаги.

\begin{enumerate}
    \item Профильный график. Субъект $\times$ условие (линии по испытуемым).
    \item Ошибка. Сравнить однофакторный ANOVA на $36$ точках
          без блока субъект с repeated-measures моделью (\texttt{AnovaRM} / mixed LM с случайным intercept).
    \item Проверить сферичность (Mauchly). При нарушении~--- поправка
          Greenhouse--Geisser.
    \item Апостериорный попарно (с поправкой) или контрасты доза и None.
    \item Интерпретация. Растёт ли время реакции с дозой?
          Почему $36$ наблюдений нельзя считать независимыми.
          Связь с \S~\ref{sec:anova} о дисперсионном анализе с повторными измерениями.
\end{enumerate}

\section{Контрольные вопросы по теме}

\begin{enumerate}
    \item Чем корреляция Пирсона отличается от Спирмена?
          Приведите пример нелинейной монотонной связи, где $\hat r_{\text{Пирсона}}\approx 0$, а $\hat r_{\text{Спирмена}}\neq 0$.
    \item Для таблицы образование--вера укажите число степеней свободы
          критерия $\chi^2$ и интерпретируйте результат в словах.
    \item Запишите разложение $Y_{ijk}=\mu+\alpha_i+\beta_j+(\alpha\beta)_{ij}+\varepsilon_{ijk}$
          для ToothGrowth. Какие факторы соответствуют \texttt{supp} и \texttt{dose}?
    \item Почему для гипертонии уместен отдельный ANOVA
          по каждому уровню \texttt{drug}?
    \item Объясните структуру времени реакции как repeated measures.
          Какой F-тест проверяет эффект условия?
    \item Сравните ANOVA и парный $t$-критерий при двух группах.
          Когда результаты совпадают?
    \item Почему корреляция между числом пожарных и ущербом от огня
          не доказывает, что пожарные «вызывают» ущерб?
    \item Вычислите $V$ Крамера для таблицы образование--вера
          при $\chi^2\approx 76{,}1$, $n$ известном из данных.
    \item Объясните парадокс Симпсона на примере лечения по полу.
    \item Запишите разложение $\mathrm{SST}=\mathrm{SSW}+\mathrm{SSB}$
          и интерпретируйте $\eta^2$ для ToothGrowth.
    \item Когда предпочтителен Крускал--Уоллис вместо F-критерия?
    \item Чем HSD Тьюки отличается от попарных $t$-критериев без поправки?
    \item Зачем при repeated measures проверять сферичность?
\end{enumerate}
